\documentclass[10.5pt,a4paper]{article}

% ============================================================
% PACKAGES
% ============================================================

\usepackage[a4paper,margin=0.62in]{geometry}
\usepackage{enumitem}
\usepackage{titlesec}
\usepackage[hidelinks]{hyperref}
\usepackage{microtype}

% ============================================================
% GENERAL FORMATTING
% ============================================================

\setlength{\parindent}{0pt}
\setlength{\parskip}{0pt}

\pagestyle{empty}

\setlist[itemize]{
    leftmargin=1.35em,
    itemsep=2pt,
    topsep=2pt,
    parsep=0pt,
    partopsep=0pt
}

% ============================================================
% SECTION FORMATTING
% ============================================================

\titleformat{\section}
  {\large\bfseries}
  {}
  {0pt}
  {}
  [\vspace{2pt}\hrule\vspace{4pt}]

\titlespacing*{\section}
  {0pt}{8pt}{3pt}

% ============================================================
% DOCUMENT
% ============================================================

\begin{document}

% ============================================================
% HEADER
% ============================================================

\begin{center}

    {\LARGE \textbf{MARRYUM ZIA}}\\[4pt]

    \textit{Physics \textbar{} Medical Physics \textbar{} Nuclear Medicine Imaging \textbar{} Artificial Intelligence}\\[5pt]

    \href{mailto:ziamarryum@gmail.com}{ziamarryum@gmail.com}
    \textbar{}
    +92 335 7500508
    \textbar{}
    Faisalabad, Pakistan\\[3pt]

    \href{https://github.com/Marryum20}{GitHub}
    \textbar{}
    \href{https://scholar.google.com/citations?view_op=new_articles&hl=en&imq=Marryum+Zia#}{Google Scholar}
    \textbar{}
    \href{https://www.linkedin.com/in/marryum-zia-3408a1243/}{LinkedIn}

\end{center}


% ============================================================
% RESEARCH PROFILE
% ============================================================

\section{Research Profile}

Physics researcher with interests in applied physics, medical physics, nuclear medicine imaging, and artificial intelligence. Experienced in computational methods, medical image analysis, deep learning, and real-world clinical imaging data. Interested in interdisciplinary research combining physics, medical imaging, and AI for clinically relevant applications.


% ============================================================
% EDUCATION
% ============================================================

\section{Education}

\textbf{University of Agriculture Faisalabad}
\hfill Faisalabad, Pakistan\\
\textit{MPhil Physics}
\hfill \textit{Oct 2024 -- Present}\\
Research focus: AI-based medical image analysis, medical physics, and
nuclear medicine imaging\\
CGPA: 3.89/4.00

\vspace{4pt}

\textbf{University of Agriculture Faisalabad}
\hfill Faisalabad, Pakistan\\
\textit{BS Physics}
\hfill \textit{Sep 2020 -- Jun 2024}\\
CGPA: 3.90/4.00


% ============================================================
% RESEARCH EXPERIENCE
% ============================================================

\section{Research Experience}

\textbf{University of Agriculture Faisalabad / PINUM Cancer Hospital}
\hfill Faisalabad, Pakistan\\
\textit{Researcher -- AI-Based Medical Image Analysis}
\hfill \textit{Jan 2026 -- Aug 2026}

\begin{itemize}

    \item Developing a deep learning framework for
    \textbf{multi-class classification of thyroid scintigraphic images}
    using a dataset of \textbf{4,639 clinically labelled images} across
    five diagnostic categories.

    \item Processing real-world clinical imaging data from
    \textbf{DICOM/PACS records}, including data extraction, anonymization,
    quality control, image cleaning, labeling, preprocessing, and dataset
    preparation.

    \item Implementing and evaluating transfer-learning architectures
    including \textbf{ConvNeXt, DenseNet, and ResNet} using PyTorch for
    automated medical image classification.

    \item Evaluating model performance using sensitivity, specificity,
    precision, F1-score, ROC-AUC, and temporally separated validation;
    \textbf{ConvNeXt-Tiny achieved approximately 91\% accuracy}
    (95\% CI: approximately \textbf{87.0--94.2\%}) on a temporally
    separated validation set of \textbf{258 images}.

    \item Applying \textbf{Grad-CAM explainability} to visualize model
    attention and improve the interpretability of AI-generated predictions
    for nuclear medicine imaging.

    \item Developing a \textbf{Streamlit-based clinical demonstration
    interface} to assess model usability and reliability and provide a
    practical framework toward future real-time clinical validation.

\end{itemize}


% ============================================================
% MEDICAL PHYSICS EXPERIENCE
% ============================================================

\section{Medical Physics Experience}

\textbf{PINUM Cancer Hospital}
\hfill Faisalabad, Pakistan\\
\textit{Medical Physicist Intern -- Nuclear Medicine Department}
\hfill \textit{Sep 2025 -- Nov 2025}

\begin{itemize}

    \item Gained hands-on experience in a clinical nuclear medicine
    environment involving \textbf{gamma cameras, SPECT/CT imaging,
    radiation instrumentation, and radiation safety}.

    \item Participated in quality control and performance evaluation
    procedures for gamma camera/SPECT imaging systems and related detector
    instrumentation.

    \item Developed practical understanding of scintigraphic imaging,
    detector operation, image acquisition, and image reconstruction in
    nuclear medicine.

    \item Gained experience with the safe handling of radioisotopes and
    radiopharmaceuticals within clinical nuclear medicine workflows.

    \item Applied concepts of nuclear physics, radiation detection,
    imaging physics, and radiation protection to real-world clinical
    applications.

\end{itemize}


% ============================================================
% PUBLICATIONS & MANUSCRIPTS
% ============================================================

\section{Publications \& Manuscripts}

\textbf{AI-Based Diagnostic Model for Thyroid Scintigraphy}\\
Manuscript submitted for publication.


% ============================================================
% SELECTED RESEARCH PROJECTS
% ============================================================

\section{Selected Research Projects}

\textbf{AI-Based Thyroid Scintigraphy Classification}
\hfill \textit{2025 -- Present}

\begin{itemize}

    \item Developed a complete medical imaging pipeline covering clinical
    data extraction, anonymization, preprocessing, deep learning,
    evaluation, explainability, and clinical demonstration.

\end{itemize}


\textbf{Technical Research Report -- Optical Materials and Instrumentation}
\hfill \textit{Jan 2024 -- May 2024}\\
\textit{University of Agriculture Faisalabad}

\begin{itemize}

    \item Reviewed the development and applications of optical materials,
    including nonlinear crystals, optical fibers, and optical coatings.

    \item Investigated applications of optical materials in detector
    windows, beam optics, and radiation-related instrumentation.

\end{itemize}


\textbf{Biometric Attendance System}
\hfill \textit{Mar 2023}\\
\textit{Undergraduate Physics Project}

\begin{itemize}

    \item Developed a fingerprint-based attendance system using Arduino
    Uno, R305 fingerprint sensor, RTC module, LCD, and SD card storage.

    \item Implemented fingerprint enrollment, verification, time-stamped
    attendance logging, and real-time system feedback.

\end{itemize}


\textbf{Lux Meter Development}
\hfill \textit{Dec 2021 -- Jan 2022}\\
\textit{Undergraduate Physics Project}

\begin{itemize}

    \item Designed and developed a digital lux meter for measuring
    illumination intensity using a light sensor and microcontroller-based
    electronic circuitry.

    \item Applied principles of electronics, sensors, measurement
    instrumentation, and experimental physics to develop and test the
    measurement system.

\end{itemize}


% ============================================================
% TECHNICAL SKILLS
% ============================================================

\section{Technical Skills}

\textbf{Programming \& Data Analysis:}
Python, NumPy, Pandas, Matplotlib, Scikit-learn

\textbf{Deep Learning \& Computer Vision:}
PyTorch, timm, Convolutional Neural Networks (CNNs),
Transfer Learning, OpenCV, Grad-CAM

\textbf{Medical Imaging \& Data:}
DICOM, pydicom, PACS data handling, medical image preprocessing,
image classification, image anonymization

\textbf{Nuclear Medicine \& Medical Physics:}
SPECT/CT, gamma cameras, scintigraphic imaging, radiation detectors,
dose calibrators, radioisotopes, radiopharmaceuticals, radiation safety,
imaging system quality control

\textbf{Software \& Research Tools:}
GitHub, VS Code, Google Colab, PyCharm, Arduino IDE, LaTeX,
Microsoft Office


% ============================================================
% HONORS & ACADEMIC ACHIEVEMENTS
% ============================================================

\section{Honors \& Academic Achievements}

\begin{itemize}

    \item Merit Scholarship recipient during BS Physics studies.

    \item Academic high achiever / class topper during undergraduate
    and postgraduate studies.

    \item Merit Scholarship recipient during MPhil Physics studies.

\end{itemize}


% ============================================================
% RESEARCH INTERESTS
% ============================================================

\section{Research Interests}

Physics \textbar{}
Medical Physics \textbar{}
Nuclear Medicine Imaging \textbar{}
Medical Image Analysis \textbar{}
Artificial Intelligence in Healthcare \textbar{}
Deep Learning \textbar{}
Computer Vision \textbar{}
Explainable AI


% ============================================================
% ENGLISH LANGUAGE PROFICIENCY
% ============================================================

\section{English Language Proficiency}

\textbf{IELTS Academic -- Overall Band 7.0}
\quad
Listening: 8.0 \textbar{}
Reading: 7.5 \textbar{}
Speaking: 7.0 \textbar{}
Writing: 6.0


\end{document}